\section{Antecedentes del proyecto}

\subsection{Situación previa}

Machu Picchu Exclusive Tours E.I.R.L. es una agencia de viajes cusqueña fundada en 2015, especializada en el sector de turismo receptivo de lujo y aventura. A pesar de contar con un producto validado y una sólida reputación basada en servicios personalizados guiados por expertos locales, la empresa operaba bajo un modelo de gestión tradicional y limitado en términos digitales.

Hasta antes de la implementación, la gestión operativa y comercial se caracterizaba por:

\begin{itemize}
    \item \textbf{Captación orgánica y manual}: los clientes, mayoritariamente extranjeros de los mercados norteamericano y europeo, contactaban a la empresa principalmente por recomendaciones directas y a través de WhatsApp, medio por el cual se realizaba todo el ciclo de venta.
    \item \textbf{Fragmentación de la información}: no se contaba con un CRM ni una base de datos estructurada. La información de los pasajeros, sus itinerarios (que varían de 2 a 20 días) y sus preferencias se gestionaba de forma dispersa en hojas de Excel, chats de WhatsApp y registros físicos.
    \item \textbf{Gestión de pagos externa}: el registro de adelantos (usualmente del 30\% al 50\%) y saldos se realizaba de manera manual, tras recibir confirmaciones de enlaces de pago externos, lo que dificultaba el control interno inmediato antes de proceder con compras no reembolsables de boletos o trenes.
    \item \textbf{Operación 100\% virtual}: al no contar con una oficina de atención física para ventas, toda la interacción y la generación de confianza dependía de videollamadas y el envío de documentos por medios digitales no centralizados.
\end{itemize}

\subsection{Problemática identificada}

El diagnóstico inicial reveló brechas críticas que obstaculizaban la escalabilidad y seguridad del negocio:

\begin{itemize}
    \item[A.] \textbf{Riesgos en la seguridad de datos}: la recepción de información sensible, como fotografías de pasaportes enviadas por WhatsApp, representaba un riesgo de ciberseguridad y privacidad, además de obligar a una depuración manual periódica para evitar el almacenamiento inseguro en dispositivos personales.
    \item[B.] \textbf{Falta de trazabilidad comercial}: la dependencia de herramientas aisladas provocaba duplicidad de tareas y pérdida de visibilidad sobre el estado real de las oportunidades de venta. No existía un historial centralizado que permitiera analizar hábitos de consumo o realizar estrategias de fidelización y remarketing.
    \item[C.] \textbf{Cuellos de botella en la gestión}: la estructura organizacional es altamente centralizada, donde la gerencia asume múltiples funciones operativas y comerciales. Sin un sistema que automatice recordatorios o el seguimiento de itinerarios, la capacidad de respuesta se veía limitada durante las temporadas de alta demanda.
    \item[D.] \textbf{Ausencia de inteligencia de negocios}: la falta de indicadores (KPIs) y reportes automatizados impedía a la empresa medir con precisión la eficacia de sus canales de captación o la rentabilidad de sus productos turísticos específicos.
\end{itemize}

